\begin{table}[htbp]
\centering
\begin{threeparttable}
\caption{Profile moderators of attacked effectivity: controlled contrasts, SEM paths, and descriptive susceptibility weights.}
\label{tab:multivariate}
\small
\setlength{\tabcolsep}{4.5pt}
\renewcommand{\arraystretch}{1.10}
\begin{adjustbox}{max width=\linewidth}
\begin{tabular}{p{0.34\linewidth}rrrrrrrrr}
\toprule
Moderator & Role & Controlled mean b & Controlled p & Ridge mean b & Weight \% & SEM mean b & SEM mean |b| & SEM min p & SEM paths \\
\midrule
Big Five Neuroticism Mean \% & core & 3.300 & 0.000 &  &  & 3.180 & 3.180 & 0.002 & 3.000 \\
Big Five Conscientiousness Mean \% & core & -0.918 & 0.293 &  &  & -0.738 & 1.059 & 0.020 & 3.000 \\
Sex Other & core & 1.819 & 0.380 & 0.063 & 2.423 & 1.527 & 3.226 & 0.058 & 3.000 \\
Big Five Openness To Experience Mean \% & core & -0.318 & 0.692 &  &  & 0.065 & 0.516 & 0.464 & 3.000 \\
Sex Female & core & 0.589 & 0.765 & -0.012 & 0.548 & 0.016 & 0.458 & 0.841 & 3.000 \\
Big Five Extraversion Mean \% & exploratory & 0.624 & 0.410 &  &  &  &  &  &  \\
Big Five Agreeableness Mean \% & exploratory & -0.186 & 0.845 &  &  &  &  &  &  \\
\bottomrule
\end{tabular}
\end{adjustbox}
\begin{tablenotes}[flushleft]
\footnotesize
\item Note. Controlled coefficients summarize moderator associations with mean attacked effectivity. SEM columns summarize the repeated-outcome path model across the attacked opinion leaves. Weight percentages come from the target-conditional regularized aggregation used to compute the post hoc susceptibility index.
\end{tablenotes}
\end{threeparttable}
\end{table}